\section*{Sets and Indices}

\begin{itemize}
    \item Feed types \(F = \{1,2,3,4,5\}\), indexed by \(f\).
    \item Nutrients \(N = \{\text{protein}, \text{minerals}, \text{vitamins}\}\). The code writes these constraints separately rather than indexing over \(N\).
\end{itemize}

\section*{Parameters}

\begin{itemize}
    \item \(c_f\): price of feed \(f\) in Y per kg (from \texttt{Price\_Y\_per\_kg}).
    \item \(p_f\): protein contribution of feed \(f\) in g per kg (from \texttt{Protein\_g}).
    \item \(m_f\): minerals contribution of feed \(f\) in g per kg (from \texttt{Minerals\_g}).
    \item \(v_f\): vitamins contribution of feed \(f\) in mg per kg (from \texttt{Vitamins\_mg}).
    \item \(P^{\min}=700\): minimum daily protein requirement (code data: \texttt{min\_protein}).
    \item \(M^{\min}=30\): minimum daily minerals requirement (code data: \texttt{min\_minerals}).
    \item \(V^{\min}=100\): minimum daily vitamins requirement (code data: \texttt{min\_vitamins}).
    \item \(\bar M = 100000\): big-\(M\) constant used in the mutual exclusion logic (code: \texttt{M}).
\end{itemize}

Using the CSV data,
\[
(c_1,c_2,c_3,c_4,c_5)=(0.2,0.7,0.4,0.3,0.8),
\]
\[
(p_1,p_2,p_3,p_4,p_5)=(3,2,1,6,18),
\]
\[
(m_1,m_2,m_3,m_4,m_5)=(1,0.5,0.2,2,0.5),
\]
\[
(v_1,v_2,v_3,v_4,v_5)=(0.5,1,0.2,2,0.8).
\]

\section*{Decision Variables}

\begin{itemize}
    \item \(x_f\) (code: \texttt{x[f]}) for \(f \in F\): nonnegative continuous amount of feed \(f\) selected.
    \item \(y_4\) (code: \texttt{y4}): binary indicator used to allow positive use of Feed 4.
    \item \(y_5\) (code: \texttt{y5}): binary indicator used to allow positive use of Feed 5.
\end{itemize}

\section*{Objective}

Minimize total feed cost:
\[
\min \sum_{f \in F} c_f x_f.
\]

Explicitly,
\[
\min\; 0.2x_1 + 0.7x_2 + 0.4x_3 + 0.3x_4 + 0.8x_5.
\]

\section*{Constraints}

Nutrient requirements:
\[
\sum_{f \in F} p_f x_f \ge P^{\min},
\]
\[
\sum_{f \in F} m_f x_f \ge M^{\min},
\]
\[
\sum_{f \in F} v_f x_f \ge V^{\min}.
\]

Explicitly,
\[
3x_1 + 2x_2 + x_3 + 6x_4 + 18x_5 \ge 700,
\]
\[
x_1 + 0.5x_2 + 0.2x_3 + 2x_4 + 0.5x_5 \ge 30,
\]
\[
0.5x_1 + x_2 + 0.2x_3 + 2x_4 + 0.8x_5 \ge 100.
\]

Mutual exclusion between Feed 4 and Feed 5:
\[
x_4 \le \bar M\, y_4,
\]
\[
x_5 \le \bar M\, y_5,
\]
\[
y_4 + y_5 \le 1.
\]

Variable domains:
\[
x_f \ge 0 \qquad \forall f \in F,
\]
\[
y_4,y_5 \in \{0,1\}.
\]

\section*{Notes on Implementation}

The implemented model in \texttt{current\_heuristic.py} is a mixed-integer linear program solved with \texttt{GUROBI\_CMD}. Feed amounts \(x_f\) are continuous and nonnegative. The incompatibility between Feed 4 and Feed 5 is modeled exactly as in the code through two binary variables and a big-\(M\) construction with \(\bar M=100000\); this formulation permits Feed 4 alone, Feed 5 alone, or neither, but not both simultaneously.

No upper bounds are imposed on feed amounts other than nonnegativity and the big-\(M\) links for Feeds 4 and 5. The nutrient constraints are modeled as lower bounds only, so oversupply is allowed if it reduces total cost.
